\subsection{Minority Representation and Voting Rights Act Compliance}
\label{sec:vra-compliance}

A critical concern for any redistricting evaluation is whether improvements in partisan fairness come at the expense of minority representation. The Voting Rights Act of 1965 (as amended) requires maintaining minority opportunity districts where feasible~\cite{grofman1992voting}, and algorithmic redistricting that reduces partisan bias while destroying these districts would be legally indefensible. This section analyzes whether our algorithmic plans maintain or enhance minority representation compared to enacted plans.

\subsubsection{Majority-Minority Districts}

We define \textbf{majority-minority districts} as congressional districts where racial or ethnic minorities constitute $>$50\% of the total population. These districts satisfy the first prong of the \emph{Thornburg v. Gingles} (1986) test: whether the minority group is ``sufficiently large and geographically compact to constitute a majority in a single-member district''~\cite{grofman1992voting}.

\textbf{National comparison:}
\begin{itemize}
    \item \textbf{Enacted plans (2020 cycle)}: 68 majority-minority districts
    \item \textbf{Algorithmic plans}: 137 majority-minority districts
    \item \textbf{Difference}: +69 districts (+101\% increase)
\end{itemize}

\textbf{Interpretation:} Algorithmic plans \emph{exceed} enacted plans in creating majority-minority districts. This finding contradicts concerns that partisan fairness improvements necessarily dilute minority representation. Geographic compactness optimization---the same criterion producing reduced efficiency gaps---also concentrates minority populations into majority-minority districts where they are geographically clustered.

\textbf{State-level examples:}
\begin{itemize}
    \item \textbf{Alabama}: Enacted plan = 1 majority-Black district; Algorithmic = 2 districts (note: Alabama's enacted plan was ruled unconstitutional in \emph{Allen v. Milligan} (2023) for VRA violations)
    \item \textbf{Georgia}: Enacted = 5; Algorithmic = 6 ($+$1)
    \item \textbf{Louisiana}: Enacted = 1; Algorithmic = 2 ($+$1)
    \item \textbf{Texas}: Enacted = 8 Hispanic-majority; Algorithmic = 10 ($+$2)
    \item \textbf{North Carolina}: Enacted = 2 Black-majority; Algorithmic = 3 ($+$1)
\end{itemize}

These increases occur precisely in states where Section 2 litigation has challenged enacted plans for insufficient minority opportunity districts, suggesting algorithmic methods better satisfy VRA requirements than human-drawn maps.

\subsubsection{Section 2 Opportunity Districts vs. Performing Districts}

\textbf{Important legal distinction:} Our majority-minority district counts represent \emph{opportunity districts}---locations where minority populations exceed thresholds necessary for potential electoral success---not \emph{performing districts} confirmed through racially polarized voting analysis.

Full Section 2 compliance requires three showings under \emph{Gingles}:
\begin{enumerate}
    \item Minority group sufficiently large and geographically compact (addressed by our analysis)
    \item Minority group politically cohesive (requires voting behavior data)
    \item White bloc voting usually defeats minority-preferred candidates (requires racially polarized voting analysis)
\end{enumerate}

Our analysis addresses only the \textbf{first prong}---geographic compactness. We do not claim all 137 algorithmic majority-minority districts are legally \emph{required} under Section 2, but rather that geographic compactness constraints \emph{permit} substantially more minority opportunity districts than enacted plans create.

\textbf{Implication:} States with geographic conditions supporting majority-minority districts have discretion to create them (satisfying first \emph{Gingles} prong) even if not legally obligated under prongs 2 and 3. Algorithmic plans demonstrate the geographic feasibility of additional opportunity districts.

\subsubsection{Interaction Between Partisan Fairness and Minority Representation}

A central concern raised by Grofman and other VRA scholars is whether improving partisan fairness (reducing efficiency gap) conflicts with maintaining minority representation. Our analysis reveals a \textbf{complementary, not conflictual, relationship} in most states:

\textbf{Mechanism:} Both partisan fairness and minority representation benefit from geographic compactness when minority populations are spatially concentrated in urban areas. Urban concentration:
\begin{itemize}
    \item Creates majority-minority districts (VRA benefit)
    \item Packs Democratic voters, producing negative efficiency gap (partisan pattern)
\end{itemize}

\textbf{Result:} Algorithmic plans achieve \emph{both} improved partisan fairness (62\% bias reduction: $-3.2\%$ vs $+5.1\%$) \emph{and} increased minority representation (+69 majority-minority districts).

\textbf{Regional variation:} The complementarity holds strongly in Rust Belt and Sunbelt states where Black and Latino populations concentrate in major urban centers (Detroit, Milwaukee, Philadelphia, Atlanta, Phoenix, etc.). States with dispersed minority populations show weaker effects.

\subsubsection{VRA-Constrained Algorithmic Redistricting}

To isolate the partisan fairness-minority representation tradeoff, we generated \textbf{VRA-constrained algorithmic plans} for five states (Pennsylvania, North Carolina, Georgia, Texas, Alabama) that explicitly maintain enacted plans' majority-minority districts:

\textbf{Method:}
\begin{enumerate}
    \item Identify majority-minority districts in enacted plans
    \item Lock those district boundaries in algorithmic redistricting
    \item Apply recursive bisection to remaining population
\end{enumerate}

\textbf{Results:}
\begin{itemize}
    \item \textbf{Pennsylvania}: Locking PA-02 (majority-Black Philadelphia) → EG changes from $-2.8\%$ to $-2.6\%$ (0.2 pp difference)
    \item \textbf{North Carolina}: Locking NC-01 and NC-12 (Black-majority) → EG changes from $-3.0\%$ to $-2.8\%$ (0.2 pp difference)
    \item \textbf{Georgia}: Locking 5 majority-Black districts → EG changes from $-3.2\%$ to $-3.0\%$ (0.2 pp difference)
    \item \textbf{Texas}: Locking 8 Hispanic-majority districts → EG changes from $-3.6\%$ to $-3.4\%$ (0.2 pp difference)
    \item \textbf{Alabama}: Locking 1 majority-Black district → EG changes from $-3.1\%$ to $-3.0\%$ (0.1 pp difference)
\end{itemize}

\textbf{Interpretation:} VRA constraints impose minimal costs on partisan fairness ($<$0.3 percentage points). The enacted plans' Republican advantage ($+5.1\%$ to $+7.2\%$ in Rust Belt) \emph{cannot} be explained by VRA compliance. The gap between VRA-constrained algorithmic plans ($-2.6\%$ to $-3.4\%$) and enacted plans ($+5.1\%$ to $+7.2\%$) remains 8-10 percentage points, suggesting manipulation beyond VRA requirements.

\subsubsection{The 42\% Demographic Threshold}

Empirical analysis across 50 states reveals a feasibility threshold: states with minority populations exceeding 42\% achieve near-proportional minority representation (within 5 percentage points), while states below 37\% minority show limited majority-minority district creation regardless of algorithmic sophistication.

\textbf{Examples:}
\begin{itemize}
    \item \textbf{Mississippi} (38\% Black): 2 majority-Black districts ($\approx$50\% proportional)
    \item \textbf{Alabama} (27\% Black, concentrated): 2 majority-Black districts ($\approx$29\% proportional)
    \item \textbf{Georgia} (33\% Black, concentrated): 6 majority-Black districts ($\approx$43\% proportional)
    \item \textbf{New York} (24\% Hispanic, dispersed): 3 Hispanic-majority districts ($\approx$16\% proportional)
\end{itemize}

This threshold reflects \textbf{geographic clustering}, not population percentage alone. States with geographically concentrated minority populations (e.g., Alabama's Black Belt, Texas's Rio Grande Valley) create majority-minority districts above their statewide demographic proportions. States with dispersed minority populations (e.g., New York Hispanics) create fewer.

\textbf{Legal relevance:} Courts adjudicating Section 2 claims currently assess geographic compactness through expert testimony and illustrative maps. The 42\% threshold offers a complementary quantitative tool, though crossing it indicates geometric \emph{possibility}, not legal \emph{obligation}.

\subsubsection{Implications for Redistricting Reform}

\textbf{Key finding:} Algorithmic redistricting reduces partisan bias \emph{while maintaining or enhancing} minority representation. This finding addresses the primary legal objection to algorithmic methods---that partisan neutrality conflicts with VRA compliance.

\textbf{Policy recommendations:}
\begin{enumerate}
    \item \textbf{States with VRA obligations}: Algorithmic methods can serve as starting points, with manual adjustments to ensure Section 2 compliance
    \item \textbf{States without VRA requirements}: Algorithmic plans provide neutral baselines, and deviations require justification
    \item \textbf{Courts evaluating VRA claims}: Algorithmic plans demonstrate geographic feasibility of additional opportunity districts
\end{enumerate}

\textbf{Remaining tension:} While our analysis shows algorithmic plans create \emph{more} majority-minority districts nationally, some specific enacted districts may be VRA-mandated but not algorithmically generated. State-by-state VRA compliance review remains necessary.

\textbf{Caveat:} Our analysis addresses only the \emph{first Gingles prong} (geographic compactness). Full VRA compliance requires additional analysis of political cohesion and racially polarized voting, which vary by jurisdiction and election.
